\documentclass[10pt,openany]{book}
\tracinglostchars=2

\makeatletter
\def\input@path{{ai almanach v3/}{ai almanach/}{./}}
\makeatother

\usepackage[USenglish]{babel}
\usepackage{fontspec}
\IfFontExistsTF{Libertinus Serif}{
  \setmainfont{Libertinus Serif}
  \setsansfont{Libertinus Sans}
}{
  \setmainfont{TeX Gyre Pagella}
  \setsansfont{TeX Gyre Heros}
}
\IfFontExistsTF{Inconsolata}{
  \setmonofont{Inconsolata}[Scale=MatchLowercase]
}{
  \setmonofont{Latin Modern Mono}
}

\usepackage[
  a4paper,
  inner=28mm,
  outer=24mm,
  top=25mm,
  bottom=27mm,
  headheight=20pt,
  headsep=8mm,
  footskip=12mm
]{geometry}
\usepackage[dvipsnames,table]{xcolor}
\definecolor{AlmanachInk}{HTML}{172033}
\definecolor{AlmanachTeal}{HTML}{167D82}
\definecolor{AlmanachOchre}{HTML}{B56D28}
\definecolor{AlmanachMist}{HTML}{EDF3F3}
\definecolor{AlmanachRed}{HTML}{9C3D35}
\color{AlmanachInk}

\usepackage{microtype}
\usepackage{setspace}
\setstretch{1.055}
\setlength{\parindent}{0pt}
\setlength{\parskip}{0.54\baselineskip}
\raggedbottom

\usepackage{titlesec}
\titleformat{\chapter}[display]
  {\normalfont\sffamily\bfseries\color{AlmanachInk}}
  {\color{AlmanachTeal}\fontsize{13}{15}\selectfont
   CHAPTER \thechapter}
  {0.45em}
  {\raggedright\fontsize{27}{31}\selectfont}
\titlespacing*{\chapter}{0pt}{-12pt}{24pt}
\titleformat{\section}
  {\Large\sffamily\bfseries\color{AlmanachInk}}
  {\color{AlmanachTeal}\thesection}{0.65em}{}
\titleformat{\subsection}
  {\large\sffamily\bfseries\color{AlmanachInk}}
  {\color{AlmanachTeal}\thesubsection}{0.6em}{}
\titleformat{\subsubsection}
  {\normalsize\sffamily\bfseries\color{AlmanachInk}}
  {\color{AlmanachTeal}\thesubsubsection}{0.55em}{}

\usepackage{fancyhdr}
\pagestyle{fancy}
\fancyhf{}
\fancyhead[LE]{\sffamily\small\thepage\quad AI Almanach}
\fancyhead[RO]{\sffamily\small\nouppercase{\leftmark}\quad\thepage}
\renewcommand{\headrulewidth}{0.3pt}
\renewcommand{\headrule}{\hbox to\headwidth{%
  \color{AlmanachTeal}\leaders\hrule height \headrulewidth\hfill}}

\usepackage{graphicx}
\usepackage{svg}
\graphicspath{{figures/}{icons/}{logos/}{ai almanach/figures/}%
  {ai almanach/icons/}{ai almanach/logos/}}
\usepackage{tikz}
\usepackage{pgfplots}
\pgfplotsset{compat=1.18}
\usetikzlibrary{
  arrows.meta,
  angles,
  calc,
  chains,
  decorations.pathreplacing,
  fit,
  matrix,
  positioning,
  quotes,
  shapes.geometric
}
\input{shapes/flowchart_styles.tex}

\usepackage{amsmath}
\usepackage{amssymb}
\usepackage{mathtools}
\usepackage{nicematrix}
\NiceMatrixOptions{cell-space-limits=1pt}
\usepackage{siunitx}
\usepackage{tabularray}
\usepackage{booktabs}
\usepackage{caption}
\captionsetup{
  font=small,
  labelfont={bf,sf,color=AlmanachTeal},
  labelsep=period,
  justification=raggedright,
  singlelinecheck=false
}
\usepackage{listings}
\lstdefinelanguage{json}{
  basicstyle=\ttfamily\small,
  string=[s]{\"}{\"},
  stringstyle=\color{AlmanachTeal!75!black},
  comment=[l]{//},
  commentstyle=\color{gray},
  morecomment=[s]{/*}{*/},
  showstringspaces=false
}
\lstset{
  basicstyle=\ttfamily\small,
  breaklines=true,
  frame=single,
  rulecolor=\color{AlmanachMist},
  backgroundcolor=\color{AlmanachMist!45},
  showstringspaces=false
}
\usepackage{enumitem}
\setlist{nosep,leftmargin=*}
\setlist[itemize,1]{label=\textcolor{AlmanachTeal}{--}}
\usepackage{multicol}
\usepackage{xparse}

\usepackage[autostyle=true]{csquotes}
\usepackage[
  backend=biber,
  style=numeric,
  sorting=nyt,
  maxbibnames=8
]{biblatex}
\addbibresource{bibliography/references.bib}
\renewcommand{\bibfont}{\small}

\newcommand{\maincolor}{AlmanachTeal}
\newcommand{\maincolspan}{4}
\newcommand{\graphscale}{0.9}
\newcommand{\keyterm}[1]{%
  \textbf{\textcolor{AlmanachTeal!78!black}{#1}}}
\newcommand{\evidencelevel}[1]{%
  \textsf{\footnotesize\color{AlmanachTeal}Evidence: #1}}

\NewDocumentEnvironment{v3topic}{m}{%
  \par\addvspace{0.75\baselineskip}%
  \begingroup
  \IfBlankF{#1}{\paragraph{#1}\mbox{}\par\nobreak}%
}{%
  \par\endgroup\addvspace{0.35\baselineskip}%
}
\newcommand{\visualseparator}{%
  \par\smallskip\noindent\color{AlmanachTeal!65}%
  \rule{24mm}{0.6pt}\par\smallskip\color{AlmanachInk}}
\newcommand{\topicline}[1]{%
  \par\smallskip{\sffamily\bfseries\color{AlmanachTeal}#1}\par}

\NewDocumentEnvironment{didacticnote}{O{Note}}{%
  \par\addvspace{0.65\baselineskip}%
  \begin{quote}\small
  {\sffamily\bfseries\color{AlmanachTeal}#1.}\quad
}{%
  \end{quote}\addvspace{0.25\baselineskip}%
}
\NewDocumentEnvironment{plainlanguage}{O{In plain language}}{%
  \begin{didacticnote}[#1]}{\end{didacticnote}}
\NewDocumentEnvironment{learningobjectives}{}{%
  \begin{didacticnote}[Learning objectives]}{\end{didacticnote}}
\NewDocumentEnvironment{engineernote}{O{Engineer's note}}{%
  \begin{didacticnote}[#1]}{\end{didacticnote}}
\NewDocumentEnvironment{workedexample}{O{Worked example}}{%
  \begin{didacticnote}[#1]}{\end{didacticnote}}
\NewDocumentEnvironment{checkpoint}{O{Check your understanding}}{%
  \begin{didacticnote}[#1]}{\end{didacticnote}}
\NewDocumentEnvironment{practicallab}{O{Practical lab}}{%
  \begin{didacticnote}[#1]}{\end{didacticnote}}
\NewDocumentEnvironment{cautionbox}{O{Common trap}}{%
  \begin{didacticnote}[#1]}{\end{didacticnote}}
\NewDocumentEnvironment{v3aside}{O{Context}}{%
  \begin{didacticnote}[#1]}{\end{didacticnote}}

\NewDocumentCommand{\chapterlead}{m m m m m m}{%
  \noindent
  \begin{minipage}[t]{0.50\linewidth}
    \vspace{0pt}
    {\sffamily\bfseries\color{AlmanachTeal}Essential question}\par
    {\large #1}\par
    \medskip
    This chapter moves from \emph{#2} through \emph{#3} and \emph{#4}
    to \emph{#5}. Read the diagram as a dependency map: each later step
    inherits assumptions and errors from the earlier ones.
  \end{minipage}\hfill
  \begin{minipage}[t]{0.44\linewidth}
    \vspace{0pt}\centering
    \begin{tikzpicture}[
      roadmap/.style={
        draw=AlmanachTeal,
        rounded corners=1mm,
        fill=AlmanachMist,
        text width=31mm,
        minimum height=7mm,
        align=center,
        font=\sffamily\small},
      link/.style={-{Stealth[length=2mm]},AlmanachOchre,thick},
      node distance=3.7mm]
      \node[roadmap] (a) {#2};
      \node[roadmap,below=of a] (b) {#3};
      \node[roadmap,below=of b] (c) {#4};
      \node[roadmap,below=of c] (d) {#5};
      \draw[link] (a) -- (b);
      \draw[link] (b) -- (c);
      \draw[link] (c) -- (d);
    \end{tikzpicture}
    \captionof{figure}{Chapter roadmap from #2 to #5.}
    \label{fig:roadmap-#6}
  \end{minipage}
  \par\vspace{1.2\baselineskip}
}

\NewDocumentCommand{\chapterreview}{m m m m}{%
  \section{Review and transfer}
  \begin{description}[style=nextline]
    \item[Reconstruct]
      Draw the path from \emph{#1} to \emph{#4} from memory. Add the
      representation, transformation, and assumption at each boundary.
    \item[Diagnose]
      Name one plausible failure in #1, #2, #3, and #4. For each one,
      identify the earliest observable signal that distinguishes it.
    \item[Verify]
      Select one equation or algorithm. Check dimensions, calculate a toy
      case, compare with a baseline, and state what the check cannot prove.
    \item[Transfer]
      Move the chapter's reasoning chain to a new domain. State which
      assumptions survive, which break, and what new evidence is required.
  \end{description}
}

\setcounter{secnumdepth}{3}
\setcounter{tocdepth}{2}

\colorlet{citecolor}{AlmanachTeal!70!black}
\colorlet{linkcolor}{AlmanachTeal!70!black}
\colorlet{urlcolor}{AlmanachOchre!85!black}
\usepackage[
  bookmarks=true,
  breaklinks=true,
  colorlinks=true,
  citecolor=citecolor,
  linkcolor=linkcolor,
  urlcolor=urlcolor,
  pdfborder={0 0 0}
]{hyperref}
\usepackage[capitalise,nameinlink]{cleveref}
\hypersetup{
  pdftitle={AI Almanach, third edition},
  pdfauthor={Bernhard Gerlach},
  pdfsubject={AI, data science, and intelligent systems for engineers},
  pdfkeywords={AI, machine learning, data science, engineering, LLMs}
}
